% !PRINTED 641
\weberpage{641}

The quotients of integral functions in $Z$ are the functionals in $\overline{Z}$. Every functional $\overline{A}$ in $\overline{Z}$ can be put into the form:
\weq{1}{\overline{A} = \frac{\Phi_1}{\Phi_2} = a\frac{E_1}{E_2} = aE,}
where $\Phi_1$, $\Phi_2$, $E_1$, $E_2$ are integral functions in $Z$, among them the units $E_1$, $E_2$, and $E$ is a unit in $\overline{Z}$ which is represented as a fractional function. $a$ is the quotient of the divisors of $\Phi_1$ and $\Phi_2$, hence an integral or fractional rational function of $z$.

The functions $a$ and $E$ in formula (1) are completely determined by $\overline{A}$, apart from \gesperrt{constant factors}, which here play the role of the \gesperrt{numerical units}. We shall call $a$ the ``\gesperrt{Absolute}'' of $\overline{A}$. The following theorem then holds:

1. \gesperrt{The Absolute of a product is equal to the product of the Absolutes of the factors,}
and we define:

2. \gesperrt{A rational functional is called integral if its Absolute is an integral function of $z$.}

From 1. it follows that the product of two integral functionals in $\overline{Z}$ is again an integral functional. The same follows also for the sum and difference of two integral functionals $\overline{A}_1$, $\overline{A}_2$ in $\overline{Z}$. For if we set:
\[
\overline{A}_1 = a_1\frac{E_1}{E}, \quad \overline{A}_2 = \frac{a_2E_2}{E},
\]
\[
\overline{A}_1 \pm \overline{A}_2 = \frac{a_1E_1 \pm a_2E_2}{E},
\]
where $a_1$, $a_2$ are integral rational functions of $z$, and $E_1$, $E_2$, $E$ integral units, then the Absolute of $\overline{A}_1 \pm \overline{A}_2$ is the divisor of the integral function $a_1E_1 \pm a_2E_2$, hence likewise an integral function of $z$. Thus the theorem holds:

3. \gesperrt{Sum, difference and product of two integral functionals in $\overline{Z}$ are again integral rational functionals.}

By the definition, all units and their reciprocals are to be designated as \gesperrt{integral}. If the Absolute of an integral rational functional $\overline{A}$ is linear ($= z - c$), then $\overline{A}$ is called a \gesperrt{rational prime functional} or a \gesperrt{prime functional} in $\overline{Z}$.
